% !TEX root = main.tex
% ======================================================================
% Beamer template for lectures (16:9 handout version)
% MIT © 2025 Magnus Simonsen
%
% This template provides:
%   - Clean 16:9 Beamer slides optimized for handouts.
%   - Color macros and tcolorboxes for structure and emphasis.
%   - Optional 2-up or 4-up handout printing layouts.
% ======================================================================

\documentclass[11pt, aspectratio=169, handout]{beamer}

% -------------------------------
% Encoding, language, typography
% -------------------------------
\usepackage[utf8]{inputenc}   % Allows typing UTF-8 characters directly
\usepackage[T1]{fontenc}      % Correct font encoding for European languages
\usepackage[english]{babel}   % Sets document language for hyphenation and formatting
\usepackage{lmodern}          % Loads the Latin Modern font (improved default LaTeX font)
\usepackage{microtype}        % Improves character spacing and text alignment automatically
\usepackage{comicneue}        % Comic Sans refined
%\usepackage[sfdefault]{atkinson}
%\usepackage{iwona}            % Humanist sans, friendly feel
%\usepackage{kurier}           % Similar to Iwona, slightly rounder
%\usepackage{aurical}          % Handwritten cursive style
%\usepackage{opensans}         % Dyslexia-friendly, clean sans-serif
%\usepackage{cabin}            % Dyslexia-friendly, humanist sans
%\usepackage[sfdefault]{roboto}    % Dyslexia-friendly, excellent clarity
%\usepackage[sfdefault]{FiraSans}  % Dyslexia-friendly, designed for screens

% -------------------------------
% Beamer theme and navigation
% -------------------------------
%\usetheme{Copenhagen}
\setbeamercovered{transparent}
\setbeamertemplate{headline}{}                % Remove header line
\setbeamertemplate{footline}[frame number]    % Frame number only
\setbeamertemplate{navigation symbols}{}      % Hide navigation icons
\setbeamertemplate{frametitle}{%
  \nointerlineskip % This removes the default vertical space above the title text
  % ht = space above title text, dp = space below title text (inside the bar)
  \begin{beamercolorbox}[wd=\paperwidth, ht=2.1ex, dp=0.8ex, leftskip=0.3cm]{frametitle}%
    \usebeamerfont{frametitle}\insertframetitle%
  \end{beamercolorbox}%
}
% Custom bullet points
\setbeamertemplate{itemize item}{\textcolor{black}{\faAngleRight }}
\setbeamertemplate{itemize subitem}{\textcolor{black}{\faAngleRight\faAngleRight}}
\setbeamertemplate{itemize subsubitem}{\textcolor{black}{\faAngleRight\faAngleRight}}

% -------------------------------
% Commonly used packages
% -------------------------------
\usepackage{amsmath, amsfonts, amsthm, amssymb} % Core tools for math equations, fonts, and symbols
\newcommand{\norm}[1]{\left\lVert#1\right\rVert}
\newcommand{\abs}[1]{\left\lvert#1\right\rvert}
\usepackage{graphicx}                           % Allows inserting and scaling images
\usepackage{xcolor}                             % Provides advanced text and background colors
\usepackage{hyperref}                           % Creates clickable links and PDF bookmarks
\usepackage{siunitx}                            % Formats numbers and physical SI units globally
\usepackage{xspace}                             % Smartly handles spaces after your custom commands
\sisetup{detect-all, per-mode=symbol, separate-uncertainty=true}
\usepackage{pgfplots}                           % Draws high-quality data plots directly in LaTeX
\pgfplotsset{compat=1.18}
\usepackage{tikz}                               % Powerful drawing engine for vector graphics
\usepackage[most]{tcolorbox}                    % Creates the advanced colored boxes for focus/theorems
\usepackage{fontawesome5}                       % Gives access to FontAwesome icons (used in \warn / \info)
%\usepackage{minted}                            % Advanced syntax highlighting for code (requires Python setup)
\usepackage[safe]{tipa}                         % Phonetic characters ('safe' prevents math symbol crashes)

% Optional: multi-slide handout layout
\usepackage{pgfpages}                           % Combines multiple slides onto a single physical printed page
% \pgfpagesuselayout{2 on 1}[a4paper, border shrink=5mm]
% \pgfpagesuselayout{4 on 1}[a4paper, border shrink=5mm]

% -------------------------------
% Hyperlink settings
% -------------------------------
\hypersetup{
  pdftitle={Title of the Lecture Series},
  pdfauthor={The Lecturer's Name},
  colorlinks=true,
  linkcolor=blue,
  urlcolor=cyan,
  citecolor=magenta
}

% -------------------------------
% Custom colors
% -------------------------------
\colorlet{myblue}{black!45!blue!80}
\colorlet{myred}{black!45!red!80}
\colorlet{mygreen}{black!45!green!80}
\colorlet{mycyan}{black!25!cyan}
\colorlet{mygray}{black!75}

% -------------------------------
% Frame color shortcuts
% -------------------------------
\newcommand{\blueheader}{%
  \setbeamercolor{frametitle}{bg=myblue, fg=white}%
}
\newcommand{\redheader}{%
  \setbeamercolor{frametitle}{bg=myred, fg=white}%
}
\newcommand{\greenheader}{%
  \setbeamercolor{frametitle}{bg=mygreen, fg=white}%
}
\newcommand{\cyanheader}{%
  \setbeamercolor{frametitle}{bg=mycyan, fg=white}%
}
\newcommand{\grayheader}{%
  \setbeamercolor{frametitle}{bg=mygray, fg=white}%
}

% -------------------------------
% Warning and info macros
% -------------------------------
\newcommand{\warn}{%
  \textcolor{red}{\faExclamationTriangle\ \textbf{Not part of the syllabus}}\xspace%
}
\newcommand{\warnCustomMsg}[2][red]{%
  \textcolor{#1}{\faExclamationTriangle\ \textbf{#2}}\xspace%
}
\newcommand{\info}{%
  \textcolor{blue}{\faInfoCircle\ \textbf{For your information}}\xspace%
}
\newcommand{\infoCustomMsg}[2][blue]{%
  \textcolor{#1}{\faInfoCircle\ \textbf{#2}}\xspace%
}

% -------------------------------
% tcolorbox environments (color-based names)
% Usage: \begin{blue}{Title}...\end{blue}
% -------------------------------
\tcbset{
  colback=white,
  fonttitle=\bfseries,
  breakable,
  enhanced,
  top=2pt,bottom=2pt,left=2pt,right=2pt,
  boxsep=3pt
}

\newtcbtheorem{bluebox}{Definition}%
{colframe=myblue}{bluebox}

\newtcbtheorem[use counter from=bluebox]{redbox}{Theorem}%
{colframe=myred}{redbox}

\newtcbtheorem[use counter from=bluebox]{greenbox}{Example}%
{colframe=mygreen}{greenbox}

\newtcbtheorem[use counter from=bluebox]{cyanbox}{Exercise}%
{colframe=mycyan}{cyanbox}

\newtcbtheorem[use counter from=bluebox]{graybox}{Explore}%
{colframe=mygray}{graybox}

% Optionally number theorems on slides:
% \setbeamertemplate{theorems}[numbered]

% -------------------------------
% Focus boxes with blue, red, green, cyan, gray background and no frame and no title with a simple command.
% We use a numbers to lighten the basic colors for each color seperate here in the definition of the focus boxes.
\newtcolorbox{bluefocus}{colback=myblue!15, colframe=white}
\newtcolorbox{redfocus}{colback=myred!15, colframe=white}
\newtcolorbox{greenfocus}{colback=mygreen!15, colframe=white}
\newtcolorbox{cyanfocus}{colback=mycyan!15, colframe=white}
\newtcolorbox{grayfocus}{colback=mygray!15, colframe=white}
